\chapter{Materials and Methods}
\label{chap:methods}

\section{Study Design}
\label{sec:study-design}
\writingnote{Present one APTOS-based study with exploratory and final evaluation
stages. Identify which artifacts belong to each stage. Keep the main narrative
compact; detailed audit and split-construction rules belong in
Appendix~\ref{app:data-protocol}.}

\figureplaceholder{Overview of the study and final evaluation workflow.}{fig:workflow}{%
Draw data preparation and grouped splitting, training on train only, checkpoint
selection using validation, and evaluation of the frozen model on test.
Show Grad-CAM as an analysis branch, not an input label or a training target.}

\section{Dataset and Target Definition}
\label{sec:dataset}
\writingnote{Document the actual APTOS distribution used, its source, original
image count, and label mapping: grade 0 becomes No DR; grades 1--4 become DR
Present. Distinguish the received CSV partitions from the original competition's
unlabelled test set. Preserve a citation to the dataset and its provenance.}

\section{Data Preparation and Partitioning}
\label{sec:data-preparation}
\writingnote{Describe the final preparation briefly and accurately: from 3,662
records, exclude eight records with conflicting binary labels and 124 redundant
exact copies, retaining 3,530 images. Keep detected similarity groups together
in the final split: 2,822 train, 354 validation, and 354 test. Explain that
image-content checks do not establish patient independence. Put hashing,
grouping, and fold-assignment details in Appendix~\ref{app:data-protocol}.}

\tableplaceholder{Class distribution in the final binary partitions.}{tab:dataset-splits}{%
Columns: partition; No DR; DR Present; total. Use the frozen split counts listed
in MATERIALS\_MAP.md. Do not mix historical and final split sizes.}

\section{Image Processing and Augmentation}
\label{sec:image-processing}
\writingnote{Describe the final RGB resizing, ImageNet normalization, and
training-only augmentations. Optional crop and CLAHE were not used in the final
binary configuration. Describe those exploratory variants separately where
their results are reported; do not present them as part of the final pipeline.}

\section{Model Architectures}
\label{sec:architectures}
\writingnote{Summarize the exploratory custom CNN, ResNet18, DenseNet121,
EfficientNet-B0, and EfficientNet-B4. Describe the final EfficientNet-B0 and its
two-output head in more detail. State which weights were pretrained and which
were initialized for the task; do not describe fine-tuning as training from scratch.}

\figureplaceholder{Architecture of the final binary classifier.}{fig:binary-architecture}{%
Show a 224 by 224 RGB input, EfficientNet-B0 feature extractor, pooling, and the
two-output classification head. Mark the actual layer used for Grad-CAM after
checking the implementation.}

\section{Training and Model Selection}
\label{sec:training}
\writingnote{Document the final configuration and computing environment using
the saved configuration and run summary. Explain validation-only checkpoint
selection and the frozen decision rule. The final run started from ImageNet
weights, not from the historical retinal-model checkpoint.}

\tableplaceholder{Training settings for the final EfficientNet-B0 experiment.}{tab:training-settings}{%
Include initialization, input size, epochs, batch size, optimizer, learning rate,
loss, augmentations, sampling, seed, device, and checkpoint selection criterion.
Use saved settings rather than parameters from the friends' thesis examples.}

\section{Evaluation Metrics}
\label{sec:metrics}
\writingnote{Define DR Present as the positive class. Explain accuracy,
precision, sensitivity/recall, specificity, F1, ROC-AUC, and confusion counts.
Distinguish binary metrics from macro/weighted multiclass metrics. Give metric
definitions here and measured values in Chapter~\ref{chap:results}.}

\section{Grad-CAM and Error Analysis Protocol}
\label{sec:gradcam-method}
\writingnote{Specify the checkpoint, target layer, target output score,
normalization, and overlay construction. Describe analysis of all six final
test errors, including maps for both output logits. Explain any additional
example-selection rule. Do not imply lesion-mask or ophthalmologist validation.}
